% Macros specific to this dissertation. Loaded after cleancampthesis.sty, so the
% template's own colours (ctcolormain, ctcoloraccessory) are available here.
%
% The research-question texts live in this file rather than in a chapter, because
% Chapter 1 states them and Chapter 8 answers them, and the two copies must not
% drift apart. Editing the wording here changes both at once.

\usepackage{ifthen}
\usepackage{tabularx}
\usepackage{tcolorbox}
\tcbuselibrary{skins,breakable}

% cleancampthesis.sty configures titlesec for \chapter through \subsection but
% leaves \paragraph without a format, which turns any use of \paragraph into an
% error. The missing run-in format is supplied here rather than by editing the
% style file.
\titleformat{\paragraph}[runin]{\normalfont\bfseries}{}{0pt}{}
\titlespacing*{\paragraph}{0pt}{0pt}{0.6em}

% Visible marker for sections that are not yet drafted, so that an intermediate
% build shows plainly what is missing.
\newcommand{\todo}{\par\noindent
  \textcolor{ctcoloraccessory}{\small[This section is not yet drafted.]}\par}

% Hyphenation exceptions. Author surnames that TeX breaks badly in the reference
% list; each was added after a specific overfull box in the bibliography.
\hyphenation{Al-neg-hei-mish}

% Long author lists occasionally leave a reference line a few points overfull.
% Allow the extra stretch inside the bibliography only, so body text is untouched.
\appto{\bibsetup}{\emergencystretch=2em}

% A measured value the author has not yet supplied (D-018). The argument names the
% quantity, so the marker doubles as the request. Every occurrence must be gone
% before the anomaly-detection chapter leaves in-progress: grep for \pending before every build.
\newcommand{\pending}[1]{\textcolor{ctcoloraccessory}{\textbf{[pending: #1]}}}

% ---------------------------------------------------------------- questions --

% Questions follow the chapter order: Android representations, modular industrial
% anomaly detection, and adversarial evaluation. RQ3 has no subquestions.
\newcommand{\rqOneText}{How can static program structure and observed network behaviour
  be jointly represented for Android malware detection and category classification, and
  how does their combination affect predictive performance compared with either source
  alone?}
\newcommand{\rqTwoText}{How can higher-order relations within static Android program
  structure be represented and learned, and to what extent do they improve malware
  detection and classification over pairwise graph representations?}
\newcommand{\rqThreeText}{How can a modular learning framework analyse network and
  process features for anomaly detection in industrial control systems?}
\newcommand{\rqFourText}{How can an explicit adversary model be translated into a
  reproducible benchmark for adversarial attacks and defences, and what does such a
  benchmark reveal about the relation between clean predictive performance and
  adversarial robustness?}

% \researchquestion{tag}{short title}{text}
% The optional first argument is a label suffix, so that \ref{rq:2} works.
\newtcolorbox{ctrqbox}[1]{%
  enhanced, breakable, sharp corners,
  colback=black!3, colframe=ctcolormain,
  boxrule=0pt, leftrule=2pt, top=6pt, bottom=6pt, left=8pt, right=8pt,
  before skip=10pt, after skip=10pt,
  fonttitle=\bfseries, title=#1, coltitle=ctcolormain,
  attach title to upper=\par\medskip}

\newcommand{\researchquestion}[4][]{%
  \ifthenelse{\equal{#1}{}}{}{\phantomsection\label{rq:#1}}%
  \begin{ctrqbox}{#2: #3}\itshape #4\end{ctrqbox}}

% ------------------------------------------------------------------- gaps --

% The four research gaps are stated once, in Section 2.7. Chapter 1 points at them
% when the questions are introduced and Chapter 8 closes them by name. Under the
% decision of 13 August 2026 the wording is NOT duplicated into a macro, because
% neither chapter reprints it; both refer back. \gapref prints the label and links
% to the statement, so a gap can never be cited by a number that does not exist.
\newcommand{\gapref}[1]{\hyperref[gap:#1]{G#1}}

% --------------------------------------------------------------- shorthands --

% System names are written as their source publications write them. Small caps were
% tried and rejected: \textsc{HGANN-Mal} renders as HGANN-Mal with the trailing
% "al" in small caps, which misreads as part of the acronym.
% The size of the retained network-flow feature set in Chapter 3. It is quoted in
% five places there, so it is fixed once here rather than retyped.
\newcommand{\netfeatures}{13}

\newcommand{\hybroid}{Hybroid}
\newcommand{\hgann}{HGANN-Mal}
\newcommand{\nadics}{NADICS}

% Project names. IUNO is an acronym and is always set in capitals, never as Iuno.
\newcommand{\iuno}{IUNO}

% Terminology fixed by the structure document. The two multiclass tasks are not the
% same task and must never be described with one phrase.
\newcommand{\familyclassification}{malware-family classification}
\newcommand{\categoryclassification}{malware-category classification}
